% Einheit 2 – Lineare Funktion f(x) = m·x + n (Nr. 10–21)
\setcounter{aufgabe}{9}
\einheitenkopf[e2][2 Lineare Funktion]{Einheit 2 von 5 · Lineare Funktion $f(x) = m \cdot x + n$}
\zweigzeile{Hier lernst du, Geraden mit $f(x) = m \cdot x + n$ zu zeichnen und $m$ und $n$ abzulesen und zu deuten · neu in diesem Jahr · P10 oft · baut auf: proportionale Funktion (Einheit 1), Brüche, negative Zahlen}

\begin{aufgabe}{Ich erkenne $m$ und $n$ in der Gleichung.\newline Unterstreiche $m$ und kreise $n$ ein. Fehlt eine Zahl, schreibe sie dazu.}
\begin{teilezwei}
\tz $f(x) = 3x + 4$ & \tz $f(x) = 5x - 2$ \\
\tz $f(x) = -4x + 1$ & \tz $f(x) = 7 - 2x$ \\
\tz $f(x) = -6$ & \\
\end{teilezwei}
\end{aufgabe}

\begin{aufgabe}{Ich erkenne ohne Zeichnen, ob eine Gerade steigt oder fällt.\newline Kreuze an.}
\begin{teile}
\teil $f(x) = 2x + 5$ \qquad \kreuz{steigt} \kreuz{fällt}
\teil $f(x) = -3x + 2$ \qquad \kreuz{steigt} \kreuz{fällt}
\teil $f(x) = 0{,}5x - 4$ \qquad \kreuz{steigt} \kreuz{fällt}
\teil $f(x) = 6 - x$ \qquad \kreuz{steigt} \kreuz{fällt}
\teil $f(x) = -2 + 3x$ \qquad \kreuz{steigt} \kreuz{fällt}
\end{teile}
\end{aufgabe}

\begin{aufgabe}{Ich kann an einer Geraden ein Steigungsdreieck einzeichnen.\newline Zeichne am markierten Punkt ein Steigungsdreieck mit einem Schritt nach rechts ein.}
\noindent\begin{minipage}[t]{0.24\linewidth}
\begin{teile}\teil\end{teile}
\begin{ksys}[xmin=-2,xmax=2,ymin=-3,ymax=3,ablesen]
\gerade{2}{-1}{}
\punkt{0}{-1}{}
\end{ksys}
\end{minipage}\hfill
\begin{minipage}[t]{0.24\linewidth}
\begin{teile}\teil\end{teile}
\begin{ksys}[xmin=-2,xmax=2,ymin=-3,ymax=3,ablesen]
\gerade{-3}{2}{}
\punkt{0}{2}{}
\end{ksys}
\end{minipage}\hfill
\begin{minipage}[t]{0.24\linewidth}
\begin{teile}\teil\end{teile}
\begin{ksys}[xmin=-2,xmax=2,ymin=-3,ymax=3,ablesen]
\gerade{1}{1}{}
\punkt{-1}{0}{}
\end{ksys}
\end{minipage}\hfill
\begin{minipage}[t]{0.24\linewidth}
\begin{teile}\teil\end{teile}
\begin{ksys}[xmin=-2,xmax=2,ymin=-3,ymax=3,ablesen]
\gerade{-2}{1}{}
\punkt{0}{1}{}
\end{ksys}
\end{minipage}
\end{aufgabe}

\verfahren{Gerade aus der Gleichung zeichnen}

\begin{aufgabe}{Ich kann zu $f(x) = m \cdot x + n$ eine Wertetabelle anlegen und die Gerade zeichnen.\newline Fülle die Tabelle aus, trage die Punkte ein und zeichne die Gerade.}
\noindent\begin{minipage}[t]{0.6\linewidth}\vspace{0pt}
\begin{teile}
\teil $f(x) = x + 2$ \quad \wertetabelle{x}{f(x)}{-1,0,1,2}
\teil $f(x) = 2x - 1$ \quad \wertetabelle{x}{f(x)}{-1,0,1,2}
\teil $f(x) = 3x - 2$ \quad \wertetabelle{x}{f(x)}{-1,0,1,2}
\end{teile}
\end{minipage}\hfill
\begin{minipage}[t]{0.38\linewidth}\vspace{0pt}
\begin{ksys}[xmin=-3,xmax=3,ymin=-5,ymax=5]
\end{ksys}
\end{minipage}
\end{aufgabe}

\begin{aufgabe}{Ich kann eine Gerade mit $n$ und Steigungsdreieck zeichnen.\newline Markiere $n$ auf der $y$-Achse, zeichne ein Steigungsdreieck und die Gerade.}
\begin{teilezwei}
\tz $f(x) = x + 3$ & \tz $f(x) = 2x + 2$ \\
\tz $f(x) = 3x + 4$ & \tz $f(x) = x + 1$ \\
\end{teilezwei}

\begin{ksys}[xmin=-4,xmax=3,ymin=-3,ymax=7]
\end{ksys}\ksysabstand\begin{ksys}[xmin=-4,xmax=3,ymin=-3,ymax=7]
\end{ksys}

\end{aufgabe}

\begin{aufgabe}{Ich kann eine Gerade mit $n$ und Steigungsdreieck zeichnen – weiter.}
\begin{teilezwei}
\tz $f(x) = -2x + 3$ & \tz $f(x) = \frac{1}{2}x + 2$ \\
\tz $f(x) = 3x - 4$ & \tz $f(x) = \frac{2}{3}x - 2$ \ (P10 2026 FOR) \\
\end{teilezwei}

\begin{ksys}[xmin=-4,xmax=4,ymin=-5,ymax=6]
\end{ksys}\ksysabstand\begin{ksys}[xmin=-4,xmax=4,ymin=-5,ymax=6]
\end{ksys}

\end{aufgabe}

\verfahren{Geraden ablesen und zuordnen}

\begin{aufgabe}{Ich kann $n$ und $m$ an einer Geraden ablesen.\newline Lies ab.}

\begin{ksys}[xmin=-4,xmax=5,ymin=-5,ymax=5,ablesen]
\gerade[2.5]{2}{-1}{g}
\gerade[-1.5]{-1}{3}{h}
\gerade[4]{0.5}{-2}{k}
\end{ksys}

\begin{teilezwei}
\tz Gerade $g$: \quad \feld{n} & \tz Gerade $h$: \quad \feld{n} \\
\tz Gerade $k$: \quad \feld{n} & \tz Gerade $g$: \quad \feld{m} \\
\tz Gerade $h$: \quad \feld{m} & \tz Gerade $k$: \quad \feld{m} \\
\end{teilezwei}
\end{aufgabe}

\begin{aufgabe}{Ich kann einer Gleichung die passende Gerade zuordnen.\newline Welche Gerade gehört zur Gleichung? Trage die Nummer der Geraden ein.}

\begin{ksys}[xmin=-4,xmax=4,ymin=-4,ymax=7,ablesen]
\gerade[-3]{2}{4}{1}
\gerade[1]{4}{2}{2}
\gerade[-1]{-2}{4}{3}
\gerade[3.5]{0}{-1}{4}
\end{ksys}\ksysabstand\begin{ksys}[xmin=-4,xmax=4,ymin=-4,ymax=7,ablesen]
\gerade[-3]{-0.5}{2}{\mathrm{I}}
\gerade[3.5]{2}{-2}{\mathrm{II}}
\gerade[3]{0}{1}{\mathrm{III}}
\end{ksys}

\begin{teilezwei}
\tz $y = 2x + 4$: \quad Gerade \leerfeld & \tz $y = -2x + 4$: \quad Gerade \leerfeld \\
\tz $y = 4x + 2$: \quad Gerade \leerfeld & \tz $y = -1$: \quad Gerade \leerfeld \\
\anweisung{Rechtes Bild: Ordne jeder Geraden ihre Gleichung zu. Drei Gleichungen bleiben übrig. (P10 2019 OS)}
\end{teilezwei}
\begin{teile}
\teil $y = 2x - 2$ \quad $y = -2x - 2$ \quad $y = -\frac{1}{2}x + 2$ \quad $y = \frac{1}{2}x + 2$ \quad $y = 1$ \quad $y = -x + 1$ \\[1mm]
I: $y = $ \leerfeld \quad II: $y = $ \leerfeld \quad III: $y = $ \leerfeld
\end{teile}
\end{aufgabe}

\verfahren{Parameter deuten}

\begin{aufgabe}{Ich kann aus $m$ und $n$ ablesen, wie eine Gerade verläuft.\newline Gegeben sind sechs Geraden. Beantworte ohne Zeichnung.}
\[ f_1(x) = 4x - 1 \qquad f_2(x) = -2x + 5 \qquad f_3(x) = 3 \qquad f_4(x) = 4x + 2 \qquad f_5(x) = -0{,}5x \qquad f_6(x) = x + 3 \]
\begin{teile}
\teil Welche Geraden steigen? \quad \leerfeld
\teil Welche Geraden fallen? \quad \leerfeld
\teil Welche Gerade verläuft parallel zur $x$-Achse? \quad \leerfeld
\teil Welche zwei Geraden sind zueinander parallel? \quad \leerfeld
\teil Welche Gerade geht durch den Ursprung $(0|0)$? \quad \leerfeld
\teil Welche zwei Geraden schneiden die $y$-Achse im selben Punkt? \quad \leerfeld
\teil Gib die Gleichung einer Geraden an, die parallel zu $f_2$ ist und die $y$-Achse bei $(0|{-1})$ schneidet. \quad $y = $ \leerfeld
\end{teile}
\end{aufgabe}

\verfahren{Prüfen, begründen, anwenden}

\begin{aufgabe}{Ich finde den Fehler beim Ablesen der Steigung.\newline Mia hat an den Steigungsdreiecken die Steigung abgelesen. Finde den Fehler, benenne ihn und korrigiere.}

\begin{ksys}[xmin=-3,xmax=3,ymin=-4,ymax=6,ablesen]
\gerade[2]{3}{-2}{u}
\steigungsdreieck{0}{-2}{3}
\gerade[-1]{-2}{4}{v}
\steigungsdreieck{0}{4}{-2}
\end{ksys}

\begin{teile}
\teil Gerade $u$: \ Mia schreibt $m = \frac{1}{3}$.
\teil Gerade $v$: \ Mia schreibt $m = 2$.
\end{teile}
\end{aufgabe}

\begin{aufgabe}{Ich kann begründen, welche Gerade steiler ist.\newline Entscheide ohne Zeichnung und begründe.}
\begin{teile}
\teil Welche Gerade ist steiler: $y = 5x - 2$ oder $y = 2x + 7$?
\teil Welche Gerade ist steiler: $y = -3x + 1$ oder $y = 2x + 1$?
\end{teile}
\end{aufgabe}

\begin{aufgabe}{Ich kann $m$ und $n$ in einer Sachsituation deuten.\newline In einer Regentonne ist Wasser. Der Hahn ist offen. Der Graph zeigt die Wassermenge $y$ in Litern nach $x$ Minuten.}

\begin{ksys}[xmin=0,xmax=16,ymin=0,ymax=70,xstep=2,ystep=10,ablesen,xlabel=$x$ in min,ylabel=$y$ in l]
\gerade[12]{-4}{60}{}
\end{ksys}

\begin{teile}
\teil Wie viel Wasser ist am Anfang in der Tonne? Welcher Parameter ist das? \quad \leerfeld
\teil Um wie viel Liter ändert sich die Wassermenge je Minute? Welcher Parameter ist das, und welches Vorzeichen hat er? \quad \leerfeld
\teil Gib die Gleichung an. \quad $y = $ \leerfeld
\end{teile}
\end{aufgabe}
